\chapter{In-depth Discussion of Existing and Proposed Solutions}
\label{ch:solutions}

\section*{Declaration of Contributions}
This chapter was jointly prepared by Parth Baranwal, Utkarsh Kumar, and Tanay Kapadia.
Parth Baranwal contributed the analysis of existing and future-oriented controls.
Utkarsh Kumar contributed MALOSS and mitigation-limitations discussion.
Tanay Kapadia contributed dependency confusion analysis and source-level detection approach.

\section{Existing Solution Families}

Current defenses against supply chain compromise can be grouped into four practical families:

\begin{enumerate}
\item \textbf{Runtime visibility and enforcement:} eBPF-based telemetry and policy enforcement
at kernel level can improve detection of anomalous process, file, and network behavior.
\item \textbf{Registry-side controls:} MFA for maintainers, typo detection, rapid advisories,
and package vetting reduce malicious publication success.
\item \textbf{Build and release integrity:} artifact signing, provenance metadata, and reproducible
build checks reduce tampering risk.
\item \textbf{Automated package analysis:} pipelines such as MALOSS combine metadata,
static, dynamic, and manual verification stages \citep{duan2021ndss}.
\end{enumerate}

\section{Dependency Confusion and Source-based Risk Assessment}

Dependency confusion exploits package name and version resolution across public and private
registries. If resolver policies prefer an untrusted public package, an attacker can hijack a
dependency path \citep{birsan2021dependencyconfusion}.

Ding \textit{et al.} proposed static source tree matching to compare intended components with
actual dependencies, enabling automated discrepancy detection even in obfuscated or injected
settings \citep{ding2024dependency}.

\section{Proposed Layered Mitigation Framework}

Based on the case study and related work, we propose a layered framework that combines
prevention, detection, and response.

\begin{algorithm}[htbp]
\caption{Layered Screening for Third-party Dependencies}
\label{alg:layered-screening}
\begin{algorithmic}[1]
\Require Dependency set $D$, policy rules $P$, trusted source map $T$
\Ensure Classified set of dependencies with risk levels
\ForAll{$d \in D$}
\State Validate source and maintainer metadata against $T$
\State Run static checks on install/build scripts and sensitive APIs
\State Trigger sandboxed dynamic checks for high-risk candidates
\State Assign risk score and enforce policy action (allow, quarantine, reject)
\EndFor
\State Publish alerts and update trust rules from analyst feedback
\end{algorithmic}
\end{algorithm}

\section{Mitigation Limits and Practical Constraints}

No single method is complete. Main constraints include:

\begin{itemize}
\item static analysis may miss runtime-triggered behavior or obfuscated logic;
\item dynamic analysis depends on execution paths and may miss dormant payloads;
\item full manual verification does not scale to very large package ecosystems;
\item delayed ecosystem response can prolong exposure windows.
\end{itemize}

Hence, effective defense is a socio-technical process that combines automation,
policy enforcement, and rapid incident communication.

\section{Chapter Summary}

The chapter demonstrates that practical risk reduction requires layered controls,
continuous monitoring, and explicit trust verification throughout the software lifecycle.
These insights lead to the consolidated conclusions and roadmap in \cref{ch:conclusion}.


